\begin{lemma}[Scheide maximum-triple fact (f)]
Let $(T,e_0,\varphi)$ be globally maximum and let $v\in V(T)$.  If $v$ is
an end of $e_0$, then
\[
 |\overline\varphi(v)|=k-d_G(v)+1\geq2;
\]
otherwise
\[
 |\overline\varphi(v)|=k-d_G(v)\geq1.
\]
\end{lemma}
\begin{proof}
The $k$-triple part of \noderef{GlobalMaximumTriple} says that $\varphi$ is
proper and leaves exactly the edge identity $e_0$ uncolored.  Properness
makes the colors on distinct incident colored edge identities at $v$
pairwise distinct.  Conversely each present color is witnessed by such an
incident edge identity by \noderef{PresentColors}.

If $v$ is incident with $e_0$, exactly $d_G(v)-1$ of its incident edge
identities are colored, by \noderef{Degree} and
\noderef{IsOneEdgeUncolored}.  Hence its present-color set has cardinality
$d_G(v)-1$.  Since \noderef{MissingColors} is the complement in the
$k$-element color set, its missing-color set has cardinality
$k-(d_G(v)-1)=k-d_G(v)+1$.  If $v$ is not incident with $e_0$, all
$d_G(v)$ incident identities are colored, giving instead
$|\overline\varphi(v)|=k-d_G(v)$.

Finally, the $k$-triple inequality in \noderef{KTriple} gives
$k\geq\Delta(G)+1$, and $d_G(v)\leq\Delta(G)$ by
\noderef{MaximumDegree}.  Thus $k-d_G(v)\geq1$; in the incident case the
additional uncolored edge raises the missing-color count to at least two.
\end{proof}
